% Options for packages loaded elsewhere
\PassOptionsToPackage{unicode}{hyperref}
\PassOptionsToPackage{hyphens}{url}
\documentclass[
]{article}
\usepackage{xcolor}
\usepackage[margin=1in]{geometry}
\usepackage{amsmath,amssymb}
\setcounter{secnumdepth}{-\maxdimen} % remove section numbering
\usepackage{iftex}
\ifPDFTeX
  \usepackage[T1]{fontenc}
  \usepackage[utf8]{inputenc}
  \usepackage{textcomp} % provide euro and other symbols
\else % if luatex or xetex
  \usepackage{unicode-math} % this also loads fontspec
  \defaultfontfeatures{Scale=MatchLowercase}
  \defaultfontfeatures[\rmfamily]{Ligatures=TeX,Scale=1}
\fi
\usepackage{lmodern}
\ifPDFTeX\else
  % xetex/luatex font selection
\fi
% Use upquote if available, for straight quotes in verbatim environments
\IfFileExists{upquote.sty}{\usepackage{upquote}}{}
\IfFileExists{microtype.sty}{% use microtype if available
  \usepackage[]{microtype}
  \UseMicrotypeSet[protrusion]{basicmath} % disable protrusion for tt fonts
}{}
\makeatletter
\@ifundefined{KOMAClassName}{% if non-KOMA class
  \IfFileExists{parskip.sty}{%
    \usepackage{parskip}
  }{% else
    \setlength{\parindent}{0pt}
    \setlength{\parskip}{6pt plus 2pt minus 1pt}}
}{% if KOMA class
  \KOMAoptions{parskip=half}}
\makeatother
\usepackage{color}
\usepackage{fancyvrb}
\newcommand{\VerbBar}{|}
\newcommand{\VERB}{\Verb[commandchars=\\\{\}]}
\DefineVerbatimEnvironment{Highlighting}{Verbatim}{commandchars=\\\{\}}
% Add ',fontsize=\small' for more characters per line
\usepackage{framed}
\definecolor{shadecolor}{RGB}{248,248,248}
\newenvironment{Shaded}{\begin{snugshade}}{\end{snugshade}}
\newcommand{\AlertTok}[1]{\textcolor[rgb]{0.94,0.16,0.16}{#1}}
\newcommand{\AnnotationTok}[1]{\textcolor[rgb]{0.56,0.35,0.01}{\textbf{\textit{#1}}}}
\newcommand{\AttributeTok}[1]{\textcolor[rgb]{0.13,0.29,0.53}{#1}}
\newcommand{\BaseNTok}[1]{\textcolor[rgb]{0.00,0.00,0.81}{#1}}
\newcommand{\BuiltInTok}[1]{#1}
\newcommand{\CharTok}[1]{\textcolor[rgb]{0.31,0.60,0.02}{#1}}
\newcommand{\CommentTok}[1]{\textcolor[rgb]{0.56,0.35,0.01}{\textit{#1}}}
\newcommand{\CommentVarTok}[1]{\textcolor[rgb]{0.56,0.35,0.01}{\textbf{\textit{#1}}}}
\newcommand{\ConstantTok}[1]{\textcolor[rgb]{0.56,0.35,0.01}{#1}}
\newcommand{\ControlFlowTok}[1]{\textcolor[rgb]{0.13,0.29,0.53}{\textbf{#1}}}
\newcommand{\DataTypeTok}[1]{\textcolor[rgb]{0.13,0.29,0.53}{#1}}
\newcommand{\DecValTok}[1]{\textcolor[rgb]{0.00,0.00,0.81}{#1}}
\newcommand{\DocumentationTok}[1]{\textcolor[rgb]{0.56,0.35,0.01}{\textbf{\textit{#1}}}}
\newcommand{\ErrorTok}[1]{\textcolor[rgb]{0.64,0.00,0.00}{\textbf{#1}}}
\newcommand{\ExtensionTok}[1]{#1}
\newcommand{\FloatTok}[1]{\textcolor[rgb]{0.00,0.00,0.81}{#1}}
\newcommand{\FunctionTok}[1]{\textcolor[rgb]{0.13,0.29,0.53}{\textbf{#1}}}
\newcommand{\ImportTok}[1]{#1}
\newcommand{\InformationTok}[1]{\textcolor[rgb]{0.56,0.35,0.01}{\textbf{\textit{#1}}}}
\newcommand{\KeywordTok}[1]{\textcolor[rgb]{0.13,0.29,0.53}{\textbf{#1}}}
\newcommand{\NormalTok}[1]{#1}
\newcommand{\OperatorTok}[1]{\textcolor[rgb]{0.81,0.36,0.00}{\textbf{#1}}}
\newcommand{\OtherTok}[1]{\textcolor[rgb]{0.56,0.35,0.01}{#1}}
\newcommand{\PreprocessorTok}[1]{\textcolor[rgb]{0.56,0.35,0.01}{\textit{#1}}}
\newcommand{\RegionMarkerTok}[1]{#1}
\newcommand{\SpecialCharTok}[1]{\textcolor[rgb]{0.81,0.36,0.00}{\textbf{#1}}}
\newcommand{\SpecialStringTok}[1]{\textcolor[rgb]{0.31,0.60,0.02}{#1}}
\newcommand{\StringTok}[1]{\textcolor[rgb]{0.31,0.60,0.02}{#1}}
\newcommand{\VariableTok}[1]{\textcolor[rgb]{0.00,0.00,0.00}{#1}}
\newcommand{\VerbatimStringTok}[1]{\textcolor[rgb]{0.31,0.60,0.02}{#1}}
\newcommand{\WarningTok}[1]{\textcolor[rgb]{0.56,0.35,0.01}{\textbf{\textit{#1}}}}
\usepackage{graphicx}
\makeatletter
\newsavebox\pandoc@box
\newcommand*\pandocbounded[1]{% scales image to fit in text height/width
  \sbox\pandoc@box{#1}%
  \Gscale@div\@tempa{\textheight}{\dimexpr\ht\pandoc@box+\dp\pandoc@box\relax}%
  \Gscale@div\@tempb{\linewidth}{\wd\pandoc@box}%
  \ifdim\@tempb\p@<\@tempa\p@\let\@tempa\@tempb\fi% select the smaller of both
  \ifdim\@tempa\p@<\p@\scalebox{\@tempa}{\usebox\pandoc@box}%
  \else\usebox{\pandoc@box}%
  \fi%
}
% Set default figure placement to htbp
\def\fps@figure{htbp}
\makeatother
\setlength{\emergencystretch}{3em} % prevent overfull lines
\providecommand{\tightlist}{%
  \setlength{\itemsep}{0pt}\setlength{\parskip}{0pt}}
\usepackage{bookmark}
\IfFileExists{xurl.sty}{\usepackage{xurl}}{} % add URL line breaks if available
\urlstyle{same}
\hypersetup{
  pdftitle={Two-sample and paired t-tests},
  hidelinks,
  pdfcreator={LaTeX via pandoc}}

\title{Two-sample and paired t-tests}
\author{}
\date{\vspace{-2.5em}}

\begin{document}
\maketitle

\textbf{Inference labs} use the five-step template: Hypothesis →
Visualise → Assumptions → Conduct → Conclude.

\subsection{Learning objectives}\label{learning-objectives}

\begin{itemize}
\tightlist
\item
  Choose between Student's and Welch's \emph{t}-test, and between paired
  and independent-samples tests.
\item
  Compute Cohen's \emph{d} and Hedges' \emph{g} as effect sizes.
\item
  Report a two-sample comparison with point estimate, CI, and effect
  size.
\end{itemize}

\subsection{Prerequisites}\label{prerequisites}

Lab 3.4.

\subsection{Background}\label{background}

The two-sample \emph{t}-test compares the means of two independent
groups. Student's version assumes equal variances across the groups;
Welch's does not. In most applied settings Welch's is the safer default
--- it behaves nearly as well as Student's when variances are actually
equal and much better when they are not. The \texttt{t.test()} default
in R is Welch.

When the same subjects are measured twice, the samples are paired and
the two measurements are correlated. Ignoring the pairing inflates the
standard error and lowers power. Pairing is handled by computing the
within-subject difference and running a one-sample \emph{t}-test on it.

Effect sizes remove units from the comparison. Cohen's \emph{d} is the
standardised mean difference: the mean difference divided by the pooled
standard deviation. Hedges' \emph{g} applies a small-sample correction.
Both are essential alongside a \emph{p}-value --- a large \emph{d} at a
small \emph{n} will often produce a non-significant \emph{p}, but the
\emph{d} is still a warning that a larger study might find a real
effect.

\subsection{Setup}\label{setup}

\begin{Shaded}
\begin{Highlighting}[]
\FunctionTok{library}\NormalTok{(tidyverse)}
\FunctionTok{library}\NormalTok{(effectsize)}
\FunctionTok{set.seed}\NormalTok{(}\DecValTok{42}\NormalTok{)}
\FunctionTok{theme\_set}\NormalTok{(}\FunctionTok{theme\_minimal}\NormalTok{(}\AttributeTok{base\_size =} \DecValTok{12}\NormalTok{))}
\end{Highlighting}
\end{Shaded}

\subsection{1. Hypothesis}\label{hypothesis}

Scenario A (independent): mean systolic blood pressure in treatment vs
placebo.

\begin{itemize}
\tightlist
\item
  H0: μ\_trt = μ\_pla, H1: μ\_trt ≠ μ\_pla, α = 0.05.
\end{itemize}

Scenario B (paired): SBP before and after a four-week intervention in
the same 30 patients.

\begin{itemize}
\tightlist
\item
  H0: mean within-subject change = 0, H1: ≠ 0.
\end{itemize}

\subsection{2. Visualise}\label{visualise}

\begin{Shaded}
\begin{Highlighting}[]
\NormalTok{trt }\OtherTok{\textless{}{-}} \FunctionTok{rnorm}\NormalTok{(n, }\AttributeTok{mean =} \DecValTok{135}\NormalTok{, }\AttributeTok{sd =} \DecValTok{12}\NormalTok{)}
\NormalTok{pla }\OtherTok{\textless{}{-}} \FunctionTok{rnorm}\NormalTok{(n, }\AttributeTok{mean =} \DecValTok{142}\NormalTok{, }\AttributeTok{sd =} \DecValTok{12}\NormalTok{)}
\NormalTok{df\_ind }\OtherTok{\textless{}{-}} \FunctionTok{tibble}\NormalTok{(}
  \AttributeTok{arm =} \FunctionTok{rep}\NormalTok{(}\FunctionTok{c}\NormalTok{(}\StringTok{"placebo"}\NormalTok{, }\StringTok{"treatment"}\NormalTok{), }\AttributeTok{each =}\NormalTok{ n),}
  \AttributeTok{sbp =} \FunctionTok{c}\NormalTok{(pla, trt)}
\NormalTok{)}
\end{Highlighting}
\end{Shaded}

\begin{Shaded}
\begin{Highlighting}[]
\NormalTok{df\_ind }\SpecialCharTok{|\textgreater{}}
  \FunctionTok{ggplot}\NormalTok{(}\FunctionTok{aes}\NormalTok{(arm, sbp, }\AttributeTok{fill =}\NormalTok{ arm)) }\SpecialCharTok{+}
  \FunctionTok{geom\_boxplot}\NormalTok{(}\AttributeTok{alpha =} \FloatTok{0.6}\NormalTok{, }\AttributeTok{colour =} \StringTok{"grey30"}\NormalTok{) }\SpecialCharTok{+}
  \FunctionTok{labs}\NormalTok{(}\AttributeTok{x =} \ConstantTok{NULL}\NormalTok{, }\AttributeTok{y =} \StringTok{"SBP (mmHg)"}\NormalTok{) }\SpecialCharTok{+}
  \FunctionTok{theme}\NormalTok{(}\AttributeTok{legend.position =} \StringTok{"none"}\NormalTok{)}
\end{Highlighting}
\end{Shaded}

\begin{Shaded}
\begin{Highlighting}[]
\NormalTok{pre  }\OtherTok{\textless{}{-}} \FunctionTok{rnorm}\NormalTok{(n2, }\DecValTok{148}\NormalTok{, }\DecValTok{10}\NormalTok{)}
\NormalTok{post }\OtherTok{\textless{}{-}}\NormalTok{ pre }\SpecialCharTok{{-}} \FunctionTok{rnorm}\NormalTok{(n2, }\DecValTok{5}\NormalTok{, }\DecValTok{7}\NormalTok{)  }\CommentTok{\# true mean drop of 5}
\NormalTok{df\_pair }\OtherTok{\textless{}{-}} \FunctionTok{tibble}\NormalTok{(}\AttributeTok{id =} \FunctionTok{seq\_len}\NormalTok{(n2), }\AttributeTok{pre =}\NormalTok{ pre, }\AttributeTok{post =}\NormalTok{ post) }\SpecialCharTok{|\textgreater{}}
  \FunctionTok{mutate}\NormalTok{(}\AttributeTok{delta =}\NormalTok{ post }\SpecialCharTok{{-}}\NormalTok{ pre)}
\end{Highlighting}
\end{Shaded}

\begin{Shaded}
\begin{Highlighting}[]
\NormalTok{df\_pair }\SpecialCharTok{|\textgreater{}}
  \FunctionTok{pivot\_longer}\NormalTok{(}\FunctionTok{c}\NormalTok{(pre, post), }\AttributeTok{names\_to =} \StringTok{"time"}\NormalTok{, }\AttributeTok{values\_to =} \StringTok{"sbp"}\NormalTok{) }\SpecialCharTok{|\textgreater{}}
  \FunctionTok{mutate}\NormalTok{(}\AttributeTok{time =} \FunctionTok{factor}\NormalTok{(time, }\AttributeTok{levels =} \FunctionTok{c}\NormalTok{(}\StringTok{"pre"}\NormalTok{, }\StringTok{"post"}\NormalTok{))) }\SpecialCharTok{|\textgreater{}}
  \FunctionTok{ggplot}\NormalTok{(}\FunctionTok{aes}\NormalTok{(time, sbp, }\AttributeTok{group =}\NormalTok{ id)) }\SpecialCharTok{+}
  \FunctionTok{geom\_line}\NormalTok{(}\AttributeTok{alpha =} \FloatTok{0.4}\NormalTok{) }\SpecialCharTok{+}
  \FunctionTok{geom\_point}\NormalTok{(}\AttributeTok{alpha =} \FloatTok{0.6}\NormalTok{) }\SpecialCharTok{+}
  \FunctionTok{labs}\NormalTok{(}\AttributeTok{x =} \ConstantTok{NULL}\NormalTok{, }\AttributeTok{y =} \StringTok{"SBP (mmHg)"}\NormalTok{)}
\end{Highlighting}
\end{Shaded}

\subsection{3. Assumptions}\label{assumptions}

Independent \emph{t}: approximate normality of the sample means; Welch's
does not require equal variances. Paired \emph{t}: approximate normality
of the within-subject differences.

\begin{Shaded}
\begin{Highlighting}[]
\NormalTok{df\_pair }\SpecialCharTok{|\textgreater{}}
  \FunctionTok{ggplot}\NormalTok{(}\FunctionTok{aes}\NormalTok{(}\AttributeTok{sample =}\NormalTok{ delta)) }\SpecialCharTok{+}
  \FunctionTok{stat\_qq}\NormalTok{() }\SpecialCharTok{+} \FunctionTok{stat\_qq\_line}\NormalTok{(}\AttributeTok{colour =} \StringTok{"firebrick"}\NormalTok{) }\SpecialCharTok{+}
  \FunctionTok{labs}\NormalTok{(}\AttributeTok{x =} \StringTok{"Theoretical"}\NormalTok{, }\AttributeTok{y =} \StringTok{"Sample (pair differences)"}\NormalTok{)}
\end{Highlighting}
\end{Shaded}

\subsection{4. Conduct}\label{conduct}

\subsubsection{Two-sample (Welch)}\label{two-sample-welch}

\begin{Shaded}
\begin{Highlighting}[]
\NormalTok{tt\_welch}
\end{Highlighting}
\end{Shaded}

\subsubsection{Two-sample Student (for
comparison)}\label{two-sample-student-for-comparison}

\begin{Shaded}
\begin{Highlighting}[]
\NormalTok{tt\_student}
\end{Highlighting}
\end{Shaded}

\subsubsection{Paired}\label{paired}

\begin{Shaded}
\begin{Highlighting}[]
\NormalTok{tt\_paired}
\end{Highlighting}
\end{Shaded}

\subsubsection{Effect size}\label{effect-size}

\begin{Shaded}
\begin{Highlighting}[]
\NormalTok{g\_ind    }\OtherTok{\textless{}{-}} \FunctionTok{hedges\_g}\NormalTok{(sbp }\SpecialCharTok{\textasciitilde{}}\NormalTok{ arm, }\AttributeTok{data =}\NormalTok{ df\_ind)}
\NormalTok{d\_paired }\OtherTok{\textless{}{-}} \FunctionTok{cohens\_d}\NormalTok{(df\_pair}\SpecialCharTok{$}\NormalTok{post, df\_pair}\SpecialCharTok{$}\NormalTok{pre, }\AttributeTok{paired =} \ConstantTok{TRUE}\NormalTok{)}
\FunctionTok{list}\NormalTok{(}\AttributeTok{d\_ind =}\NormalTok{ d\_ind, }\AttributeTok{g\_ind =}\NormalTok{ g\_ind, }\AttributeTok{d\_paired =}\NormalTok{ d\_paired)}
\end{Highlighting}
\end{Shaded}

\subsection{5. Concluding statement}\label{concluding-statement}

\begin{quote}
\textbf{Independent comparison.} Mean SBP was
\texttt{round(mean(trt),\ 1)} mmHg (SD \texttt{round(sd(trt),\ 1)}) in
the treatment arm and \texttt{round(mean(pla),\ 1)} (SD
\texttt{round(sd(pla),\ 1)}) in placebo. A Welch's \emph{t}-test gave
\emph{t}(\texttt{round(tt\_welch\$parameter,\ 1)}) =
\texttt{round(tt\_welch\$statistic,\ 2)}, \emph{p} =
\texttt{signif(tt\_welch\$p.value,\ 2)}, mean difference
\texttt{round(diff(rev(tt\_welch\$estimate)),\ 1)} mmHg (95\% CI
\texttt{round(-tt\_welch\$conf.int{[}2{]},\ 1)} to
\texttt{round(-tt\_welch\$conf.int{[}1{]},\ 1)}). Cohen's \emph{d} =
\texttt{round(d\_ind\$Cohens\_d,\ 2)} (Hedges' \emph{g} =
\texttt{round(g\_ind\$Hedges\_g,\ 2)}), a small-to-medium effect.
\end{quote}

\begin{quote}
\textbf{Paired comparison.} The within-subject change (post − pre) had
mean \texttt{round(mean(df\_pair\$delta),\ 1)} mmHg (SD
\texttt{round(sd(df\_pair\$delta),\ 1)}); paired \emph{t}-test
\emph{t}(\texttt{tt\_paired\$parameter}) =
\texttt{round(tt\_paired\$statistic,\ 2)}, \emph{p} =
\texttt{signif(tt\_paired\$p.value,\ 2)}, 95\% CI
\texttt{round(tt\_paired\$conf.int{[}1{]},\ 1)} to
\texttt{round(tt\_paired\$conf.int{[}2{]},\ 1)}.
\end{quote}

The effect size is the number that generalises. A \emph{p}-value tells
you whether to reject; \emph{d} and its CI tell you how much.

\subsection{Common pitfalls}\label{common-pitfalls}

\begin{itemize}
\tightlist
\item
  Running an independent-samples \emph{t} on paired data and losing
  power.
\item
  Reporting Student's \emph{t} in the era of Welch; there is rarely a
  good reason.
\item
  Reporting \emph{p} without \emph{d} (or an equivalent effect size).
\item
  Confusing effect size with statistical significance.
\end{itemize}

\subsection{Further reading}\label{further-reading}

\begin{itemize}
\tightlist
\item
  Delacre M et al.~(2017). \emph{Why psychologists should by default use
  Welch's t-test}.
\item
  Cohen J. \emph{Statistical Power Analysis for the Behavioral
  Sciences}.
\end{itemize}

\subsection{Session info}\label{session-info}

\begin{Shaded}
\begin{Highlighting}[]

\end{Highlighting}
\end{Shaded}

\subsection{See also --- chapter index}\label{see-also-chapter-index}

\begin{itemize}
\tightlist
\item
  \href{../index.Rmd}{Methods in this chapter}
\item
  \href{../../../ch00-find-your-method.Rmd}{Find your method (Chapter
  0)}
\end{itemize}

\end{document}
